\chapter{Conclusion}
\label{ch:conclusion}
 
\section{The Argument in Brief}
 
The \textit{Dobbs v. Jackson Women's Health Organization} decision of June 2022 posed a question the accumulated diffusion literature had not been obliged to answer: when a constitutional doctrine half a century old is vacated in a single opinion, returning to a highly salient and heavily organized policy domain in full to the states, what shape does the resulting legislative activity take? That question is precisely the empirical puzzle the received frameworks of policy diffusion were not designed to address. The question is not merely substantive. It is also diagnostic: the patterns that emerge under such conditions reveal which theories capture the underlying dynamics of state policymaking and which capture only the equilibrium arrangements to which those dynamics had adapted during the half century of federal preemption that \textit{Roe} and \textit{Casey} imposed.
 
This thesis has assembled the first systematic answer for the first full legislative year after the vacuum opened. Drawing on a corpus of 188 enacted bills across 46 states, and analyzing that corpus through computational text measures and dyadic regression, it documents a diffusion signature that neither the geographic-learning tradition nor the dominant pre-\textit{Dobbs} intuition about restrictive templating anticipated. Post-\textit{Dobbs} abortion legislation in 2023 was characterized by moderate textual borrowing across state lines---mean cross-state novelty of 0.769 with a standard deviation of 0.122, two-thirds of bills clustered between 0.70 and 0.90---a distribution consistent with widespread customization rather than wholesale copying or independent drafting. Within that baseline, the coalition that deployed standardized templates most aggressively was the protective coalition, not the restrictive one ($\beta = -0.076$, $p = 0.022$ for protective bills; $\beta = -0.017$, $p = 0.482$ for restrictive bills, under Model 3 with Bell--McCaffrey CR2 inference on $N = 175$). Textual convergence across state pairs was predicted by shared organizational infrastructure ($\beta = +0.039$, $p < 0.001$, measured as the Jaccard index of PAC donor organizations) and by shared protective coalition membership ($\beta = +0.037$, $p < 0.001$), not by geographic contiguity ($\beta = +0.003$, $p = 0.734$). The ten highest-similarity cross-state bill pairs registered zero common PAC donors, including the Nevada SB 370--Washington HB 1155 pair whose 80.7 percent cosine similarity constitutes the singular case of near-replication in the dataset.
 
The theoretical apparatus developed in Chapter 5 draws these results into a single framework. \textit{Constitutional vacuum diffusion} names the regime that emerges when federal preemption collapses abruptly and sub-national jurisdictions are compelled to legislate across a newly opened frontier within compressed timelines. \textit{Reactive template mobilization} names the mechanism that operates within that regime: pre-positioned model language, developed over years by organizations arrayed on both sides of the issue, is activated through affiliate channels in the months immediately following the shock. The coalition asymmetry the mechanism produces is not incidental to the empirical record but diagnostic of two distinct organizational architectures. Restrictive coordination in 2023 flowed through vertical template distribution from a small number of national nodes---Americans United for Life, Susan B. Anthony Pro-Life America, Alliance Defending Freedom---whose pre-\textit{Dobbs} template development was institutionally mature but whose post-shock deployment was constrained by an unsettled litigation environment in which previously drafted templates required re-engineering for the expanded policy frontier. Protective coordination, by contrast, flowed through horizontal affiliate networks whose convergence on a shared codification agenda produced the concentrated textual signature the corpus displays.
 
What follows in this concluding chapter is not a restatement of these findings but an assessment of what they contribute and where they locate the thesis within the longer trajectory of diffusion scholarship. Section 7.2 articulates the empirical, theoretical, and methodological contributions in their most concise form. Section 7.3 situates the thesis in relation to the literatures from which it departs. Section 7.4 considers its implications for post-shock policy domains beyond reproductive rights. Section 7.5 closes with the interpretive question the 2023 record poses most urgently for the field: whether the patterns documented here represent a transitional anomaly of the immediate post-\textit{Dobbs} mobilization period or the first evidence of a structural transformation in how American states make policy when federal doctrinal infrastructure collapses.
 
\section{What the Thesis Contributes}
 
\subsection{Empirical Contributions}
 
The thesis supplies the first systematic cross-state textual analysis of the post-\textit{Dobbs} abortion corpus. Prior scholarship on post-\textit{Dobbs} state legislative activity has been descriptive or case-based, documenting particular enactments without placing them on a comparative structural surface. By constructing the full enacted-bill corpus for 2023, computing pairwise textual similarity across its 188 bills, and measuring novelty as one minus each bill's maximum cross-state similarity, the analysis produces a quantitative foundation on which the diffusion dynamics of a full legislative cycle can be read directly from the statutory text.
 
The quantification of coalition asymmetry is the most consequential empirical finding. Three decades of documented model-bill production by national restrictive organizations had generated an expectation that post-shock legislative coordination would be most visible on the restrictive side. The 2023 record inverts that expectation at conventional levels of statistical significance. Protective bills exhibit significantly lower novelty than the neutral baseline; restrictive bills do not. At the dyadic level, the Both Protective coefficient is robust across specifications at $\beta = +0.037$ ($p < 0.001$), a magnitude equal to 30 percent of the dependent variable's mean (0.037/0.125), while the Both Restrictive coefficient attenuates from marginal significance in Model D2 ($\beta = +0.015$, $p < 0.10$) to null in Model D3 once organizational co-presence is controlled for ($\beta = +0.012$, $p = 0.158$). The attenuation itself is substantively informative: it is consistent with restrictive coordination operating substantially through the same organizational infrastructure that the Jaccard measure captures, while protective coordination operates through additional channels the measure does not fully register.
 
The decoupling of financial transfer from textual coordination is the second empirical contribution. Across the ten highest-similarity cross-state pairs, direct PAC contributions from abortion-focused donors to the sponsors of the converging bills approach zero. Systematic analysis of the DIME database's full contribution record confirms that abortion-focused organizations contribute negligibly to state legislative candidates in general, not merely to the sponsors of the specific bills this thesis analyzes. The dyadic regression nonetheless identifies shared PAC donor organizations as a strong and independent predictor of similarity. The reconciliation, developed at length in Section 5.6, is that the organizations whose presence the Jaccard index captures are not purchasing legislative behavior; they are supplying informational infrastructure---model language, legal memoranda, drafting assistance, expert testimony, coalition talking points---whose effects are visible in the textual record but invisible in the contribution record. The finding supplies empirical confirmation of the \textcite{hallLobbyingLegislativeSubsidy2006a} legislative subsidy framework in a policy domain where the financial-influence intuition has been especially prominent in public commentary, and it demonstrates the magnitude of error that contribution-only measures of interest-group influence can produce in domains where the informational channel dominates.
 
\subsection{Theoretical Contributions}
 
The constitutional vacuum diffusion framework developed in Chapter 5 is the thesis's principal theoretical contribution. The framework identifies a class of diffusion events characterized by four structural features: a judicially imposed rather than discursively negotiated policy window; a bilateral rather than unilateral opportunity structure that compels both coalitions to legislate simultaneously; a compressed rather than extended temporal horizon measured in months rather than years; and an advocacy infrastructure that pre-positions model language before the shock rather than producing it in response. These features distinguish constitutional vacuum diffusion from the policy-window dynamics of \textcite{kingdonAgendasAlternativesPublic2011a}, the punctuated-equilibrium dynamics of \textcite{baumgartnerLobbyingPolicyChange}, and the event-history diffusion of \textcite{berryStateLotteryAdoptions1990}. The framework's value for the broader literature is that it specifies the structural conditions under which a diffusion regime of this kind should emerge, enabling subsequent research to evaluate whether other domains in which federal preemption has collapsed---voting rights after \textit{Shelby County v. Holder} (2013), firearms regulation after \textit{New York State Rifle \& Pistol Association v. Bruen} (2022), or hypothetical reversals of other long-standing constitutional doctrines---generate comparable signatures.
 
Reactive template mobilization, developed in Section 5.3, supplies the mechanism that operates under the framework's scope conditions. Unlike learning-based diffusion, which presupposes that adopting states have time to observe the consequences of others' policy choices, and unlike conventional crisis adoption, which presupposes that templates are constructed in response to the precipitating event, reactive template mobilization presupposes that templates are pre-built and activated. The mechanism is distinct from both and has predictive implications the 2023 record confirms: convergence should concentrate on pre-existing model language rather than newly drafted text; state-to-state timing should be determined by affiliate coordination rather than observational learning; and the diffusion signature should be detectable in the statutory text within the first post-shock legislative cycle, before measurable policy outcomes emerge in the adopting states.
 
The coalition-architecture distinction---vertical restrictive templating through a small number of national nodes, horizontal protective coordination through distributed state affiliates---is the third theoretical contribution. It refines the \textcite{jansaCopyPasteLawmaking2019} copy-and-paste diffusion framework, which was developed on pre-\textit{Dobbs} evidence and predicts that the restrictive coalition should drive observable textual coordination in abortion policy, by showing that when both coalitions are activated simultaneously, their different organizational architectures produce different textual signatures. The restrictive architecture yields convergence that is less concentrated within coalition-aligned dyads; the protective architecture concentrates convergence within coalitionally aligned dyads. The architectural distinction thereby extends rather than refutes the copy-and-paste tradition, integrating it into a fuller account that specifies how coalition structure conditions the textual signature diffusion produces.
 
The reinterpretation of the professionalism--novelty relationship, developed in Section 5.5, is the thesis's fourth theoretical contribution. \textcite{squireMeasuringStateLegislative2007} and the broader state-politics literature have treated legislative capacity as enabling innovation. The 2023 record suggests instead that capacity enables adoption: professionalized legislatures exhibit lower novelty in a pattern consistent with greater capacity to exploit externally available template infrastructure, while low-capacity legislatures produce textually more original bills in a pattern consistent with restricted access to the same infrastructure. The inference about \emph{why} the relationship runs in this direction---that advocacy-supplied templates are the relevant infrastructure---is theoretically grounded but not directly established by the regression and remains, at this stage of the research program, an interpretation rather than a measured mechanism. The reinterpretation is consistent with the \textcite{jansaCopyPasteLawmaking2019} finding that professionalism correlates with template uptake in the pre-\textit{Dobbs} record, and it extends that finding to the post-shock regime. It also yields a cross-domain prediction that the thesis identifies but does not evaluate: the negative professionalism--novelty relationship should be strongest in high-infrastructure domains in which organized advocacy has produced dense template libraries, and should weaken or reverse in domains where such infrastructure is sparse.
 
The revision of the exporter--importer--reinforcer typology developed in Section 5.7 is the fifth theoretical contribution. The revised typology distinguishes exporter function, the origination of a framework that diffuses widely, from exporter status, the maintenance of high original-drafting rates across a state's legislative output. It acknowledges that the importer--exporter axis operates within coalition structures rather than across them, and it integrates the Vermont hub finding as evidence that affiliate position rather than institutional capacity determines exporter function within the horizontal protective architecture. The typology, so revised, constitutes a more precise instrument for subsequent diffusion research in domains where coalition architectures may be asymmetric and where the conventional mapping between legislative capacity and diffusion function cannot be assumed to hold.
 
\subsection{Methodological Contributions}
 
The research design contributes three methodological refinements with value beyond the immediate case. The integration of bill-level and dyadic specifications within a single analytical pipeline demonstrates the value of nested units of analysis in diffusion research. The bill-level regression identifies the institutional and coalition predictors of novelty at the level of the enacted bill; the dyadic regression identifies the structural predictors of convergence at the level of the state pair. Each specification answers questions the other cannot, and together they produce a more complete characterization of the diffusion signature than either could alone. The Bell--McCaffrey CR2 Satterthwaite inference employed throughout---implemented via \texttt{estimatr::lm\_robust}---ensures that conservative small-sample corrections appropriate to the dyadic sample size are applied uniformly, protecting against the overstatement of significance that conventional cluster-robust inference can produce in samples of this dimension \parencite{pustejovskyCRVEAnywhere2018}.
 
The Jaccard index of shared PAC donor organizations illustrates how financial data can be repurposed to measure organizational co-presence rather than financial transfer. The empirical payoff, documented in Section 4.8 and reiterated in the empirical summary above, is that the two measures dissociate sharply in the 2023 record. The methodological implication is more general: financial presence and financial transfer are distinct constructs, and diffusion research that conflates them risks systematically mismeasuring the mechanism. Subsequent studies requiring a measure of organizational infrastructure rather than financial influence can apply the construction developed here to other policy domains with comparable advocacy structures, substituting domain-appropriate donor taxonomies into the Jaccard denominator and numerator.
 
Finally, the replication infrastructure assembled for this thesis---the archival analytical files, the complete \texttt{R} codebase for both bill-level and dyadic specifications, the full $188 \times 188$ similarity matrix, and the authoritative dependent-variable construction---provides a public good for subsequent research on post-\textit{Dobbs} state abortion policy. The analytical pipeline is reproducible end-to-end from the raw legislative corpus through the final regression tables, and the documentation is sufficient for independent replication or extension. The value of such infrastructure is particularly salient for a policy domain in which the first-year record is already being cited in litigation briefs, legislative debates, and policy commentary, and in which the availability of transparent, reproducible textual measures may substantially affect the quality of downstream scholarly and practical use.
 
\section{Locating the Thesis Within Diffusion Scholarship}
 
The relationship between this thesis and the diffusion literature it draws on is neither strictly continuous nor strictly disruptive. The core insight of \textcite{walkerDiffusionInnovationsAmerican1969}---that policy spreads through structured interstate pathways rather than through simultaneous independent adoption---is confirmed by the 2023 record. What requires revision is the specification of the pathways themselves. Geographic contiguity is null in the dyadic regression ($\beta = +0.003$, $p = 0.734$); coalitional and organizational ties are strong. The underlying analytical move---modeling diffusion as a structured process with identifiable transmission channels---is preserved; the substantive content of the channels is updated to reflect the post-\textit{Dobbs} evidence. The thesis extends rather than displaces the Walker tradition, demonstrating that the analytical framework survives the transition from a geographically organized policy environment to a coalitionally organized one.
 
A parallel cumulative relationship holds with respect to \textcite{berryStateLotteryAdoptions1990}. The event-history framework's pairing of internal determinants with external influence remains analytically indispensable; the 2023 record indicates that in a constitutional vacuum regime, the external-influence term should be specified as organizational co-presence and coalition alignment rather than as regional diffusion. The framework's structure survives; its parameters are updated. Future event-history analyses of post-shock diffusion domains can inherit the general design while incorporating the revised specification the thesis supplies.
 
\textcite{jansaCopyPasteLawmaking2019}'s copy-and-paste tradition, by contrast, requires more intricate engagement. That tradition, developed on pre-\textit{Dobbs} evidence, generated an expectation of restrictive-templating dominance that the 2023 record does not confirm. The coalition asymmetry documented here is not a refutation of the copy-and-paste framework but an extension of it. The framework correctly identifies interest-group drafting as a diffusion channel and correctly predicts that the channel produces detectable textual signatures; it requires amendment only in its prediction that the signatures should be concentrated on one side of the abortion debate. The thesis's revision---that coalition architecture, not coalition identity, determines which signatures appear most prominently---is a friendly amendment to a broadly correct framework, and it carries a testable implication: where architectural asymmetries between coalitions exist, the textual signature of diffusion should track the asymmetry rather than the substantive valence of the policy under adoption.
 
The engagement with \textcite{hertel-fernandezStateCapture2019} is primarily one of empirical grounding. State Capture argues that conservative advocacy organizations have developed a systematic capacity to shape state legislative activity through model-bill distribution, affiliate coordination, and staff liaison rather than through direct campaign contributions. The 2023 abortion record provides high-resolution textual evidence that mirrors Hertel-Fernandez's case-based argument on the restrictive side while documenting a parallel---though architecturally distinct---coordination apparatus on the protective side. The thesis thereby extends the \textit{State Capture} analytical frame to a domain and a coalition the book did not extensively treat, and it supplies the text-as-data measures that the qualitative design did not develop.
 
The most substantively novel engagement is with the four-mechanism framework of \textcite{shipanMechanismsPolicyDiffusion2008a}---learning, competition, imitation, coercion. Reactive template mobilization is not straightforwardly a fifth mechanism; it is a mode of operation that can involve elements of imitation and, in attenuated form, coercion, while differing from each in its temporal structure and its dependence on pre-positioned organizational infrastructure. The contribution is not to supplement the four-mechanism typology with an additional category but to show that under constitutional vacuum conditions, the relative weight and operational structure of the four mechanisms shifts in ways the original formulation did not centrally theorize. The 2023 record does not fit cleanly into any single mechanism; it fits most closely into imitation mediated by shared organizational networks, a hybrid the four-mechanism framework contemplates without fully specifying. Subsequent diffusion research engaging seriously with post-shock legislative activity should expect to encounter such hybrids and should develop the analytical apparatus to distinguish them from the more conventional mechanisms Shipan and Volden catalogued.
 
The broader positioning that emerges from these engagements is that the thesis operates within the diffusion literature's received analytical architecture while revising its substantive specifications for a regime the literature has not previously been required to address. The geographic tradition supplies the modeling logic; the event-history tradition supplies the dependent-variable conception; the copy-and-paste tradition supplies the measurement strategy; the state-capture literature supplies the organizational framing; the four-mechanism typology supplies the mechanism vocabulary. The thesis's contribution, in aggregate, is to draw these received resources into a framework capable of characterizing a diffusion regime---constitutional vacuum diffusion---whose structural features none of the contributing traditions was designed to capture alone.
 
\section{Implications for Post-Shock Policy Domains}
 
The detailed cross-domain extensions are addressed at length in Section 6.6.4, which evaluates the framework's travel prospects for firearms regulation following \textit{Bruen}, voting rights following \textit{Shelby County}, and the hypothetical case of marriage-equality retrenchment. What warrants attention here is not the comparative case work itself but the analytical implications the thesis carries for the broader class of policy domains that may eventually experience constitutional shocks of comparable magnitude. Three points bear emphasis.
 
The availability of pre-positioned templates is the structural precondition that the reactive template mobilization mechanism requires. Domains in which organized advocacy has spent decades developing model legislation for constitutional contingencies that had not yet materialized---domains in which pre-\textit{Dobbs} abortion policy is representative rather than exceptional---should be expected to produce the compressed, text-heavy diffusion signature the 2023 record exhibits. Voting rights after \textit{Shelby County}, firearms regulation after \textit{Bruen}, and a hypothetical future retrenchment on marriage equality all satisfy this precondition to varying degrees; each involves advocacy networks whose template-production capacity developed in anticipation of constitutional change rather than in response to it. Domains lacking such pre-positioned infrastructure---emerging regulatory areas, novel technologies, or policy questions that had not been legally contested before the precipitating event---should produce a different diffusion profile in which customization is more pronounced and cross-state convergence less concentrated.
 
The coalition-architecture distinction developed in Chapter 5 should not be assumed to travel intact across policy domains. The vertical--horizontal asymmetry the 2023 abortion record exhibits is a product of the specific organizational histories of the two coalitions, not a structural feature of advocacy networks in general. Domains in which both coalitions are organized through vertically integrated national hubs---certain areas of environmental regulation and federal labor policy, for instance---should produce more symmetric diffusion signatures than the one documented here. Domains in which both coalitions are organized horizontally should produce the opposite. What the thesis contributes is not a universal claim about coalition architecture but a set of analytical tools for diagnosing the architecture of any given domain from its observed textual signature, and for predicting which coalition's convergence should be most visible in the textual record given its organizational form.
 
The measurement strategy developed here is portable to any policy domain in which enacted-bill text is available and in which organizational co-presence can be operationalized. The computational pipeline does not depend on features of abortion politics that are not shared with other domains. Its portability means that the next constitutional shock, if it produces a comparable legislative response, can be analyzed within months of its occurrence rather than the years typically required for diffusion scholarship to catch up to unfolding political events. The analytical infrastructure, pre-built, is itself a form of template; the thesis's methodological contribution is partly to demonstrate that such infrastructure can be pre-built.
 
\section{Transitional Anomaly or Structural Transformation}
 
The closing question Chapter 6 identifies as the one the 2023 record poses most urgently for the field is whether the patterns documented here represent a transitional anomaly of the immediate post-\textit{Dobbs} mobilization period or the first evidence of a structural transformation in how American states make policy after constitutional shocks. The question cannot be resolved from a single legislative year. What can be specified is what the two possibilities imply and what evidence would distinguish them.
 
The transitional-anomaly reading is that the 2023 coalition asymmetry reflects a one-time mobilization advantage produced by the particular timing of the \textit{Dobbs} decision. Protective coalitions, on this reading, had developed their codification agenda over the preceding decade in anticipation of the erosion of \textit{Roe} but had been unable to enact it against the constitutional backdrop \textit{Roe} provided; \textit{Dobbs} removed the backdrop and created an immediate legislative imperative to codify protections that had previously been assumed rather than enacted. Restrictive coalitions, correspondingly, had developed their template infrastructure against a constitutional regime that constrained permissible restrictions, and the 2023 legislative year found them re-engineering those templates for the expanded policy frontier rather than deploying pre-existing ones intact. On this reading, the coalition asymmetry should attenuate in subsequent legislative years as restrictive templates catch up to the expanded frontier and protective codification exhausts its initial agenda, producing a more symmetric diffusion signature in 2024, 2025, and beyond. The prediction is directional and testable: the longitudinal extension specified in Section 6.6.3 constitutes precisely the empirical test the question requires.
 
The structural-transformation reading is that the 2023 coalition asymmetry reflects deeper differences in how the two coalitions organize and will persist beyond the initial codification cycle. On this reading, the horizontal affiliate architecture characterizing the protective coalition is not a one-time mobilization advantage but the structural form the coalition has taken; the vertical architecture characterizing the restrictive coalition is likewise structural and historically durable. If these architectures are stable features of the respective coalitions, the textual signatures they produce should be stable as well. Protective convergence should continue to exceed restrictive convergence as both coalitions legislate against the post-\textit{Dobbs} constitutional backdrop, because the mechanism that produces convergence---affiliate coordination in the protective case, national-hub distribution in the restrictive case---operates through architecturally distinct channels on the two sides of the debate.
 
The two readings are not mutually exclusive. The 2023 coalition asymmetry may reflect both a transitional mobilization advantage on the protective side and a durable architectural difference between the coalitions, with the former dissipating as subsequent legislative cycles unfold and the latter persisting. What matters analytically is that the two components are distinguishable: a transitional advantage should attenuate, while an architectural difference should persist. Five years of post-\textit{Dobbs} legislative records, analyzed through the same computational pipeline developed here, should supply the evidence needed to distinguish the two components, and the decomposition will carry consequences for the theoretical framework this thesis has advanced.
 
The stakes of the distinction extend beyond abortion policy. If the coalition asymmetry is transitional, the inference for other constitutional-shock domains is that the coalition better positioned at the moment of the shock will exhibit the most visible textual signature, and that the signature will attenuate as the other coalition catches up. If the asymmetry is structural, the inference is that coalitions' organizational architectures produce different signatures regardless of mobilization timing, and that the textual record will consistently favor whichever coalition is better suited to distributed affiliate coordination in that domain. The interpretation of textual-convergence patterns in future post-shock domains depends sensitively on which reading prevails. The 2023 record is insufficient to adjudicate the question; it is sufficient to pose it precisely enough that subsequent research can answer it.
 
\section{Final Remarks}
 
The \textit{Dobbs} decision returned a half-century-old constitutional question to the states on a legislative timeline measured in weeks. The first year of the legislative response is now visible in its full textual detail, and the pattern it displays is not what the accumulated diffusion literature would have predicted. Moderate customization at the aggregate level, coalition asymmetry at the specification level, organizational coordination without financial transfer at the mechanism level, and a horizontal affiliate architecture on the protective side that contradicts both pre-\textit{Dobbs} expectations and popular commentary about the post-shock legislative landscape---these are the signatures the 2023 record displays, and they are consistent enough across specifications to warrant the interpretive apparatus this thesis has developed.
 
What the thesis does not claim, and what the limitations chapter has taken care to specify, is that any single 188-bill corpus can settle the large interpretive questions the diffusion literature must now confront. One legislative year is not a trend. One policy domain is not a theory. The computational and documentary evidence assembled here provides the most rigorous account currently available of the post-\textit{Dobbs} diffusion record, but the account is pressure-tested only within the year and the domain for which it was constructed. The contributions it makes to the broader literature are contingent on the patterns documented here being confirmed, revised, or refuted by the subsequent research the limitations chapter has specified.
 
That such pressure-testing will be possible is itself a feature of the research design. The analytical infrastructure is transparent; the data are public; the methods are replicable; the extensions are identified. Subsequent work, undertaken with different instruments and on different corpora, should be able to engage the thesis's claims directly rather than approximately. Whether that subsequent work confirms or revises what this thesis has documented is, appropriately, a question for the research that follows rather than for the work that concludes here. What my work establishes, within the scope conditions the design imposes, is that policy diffusion after constitutional shocks operates through channels the field has not yet fully theorized---and that the analytical tools required to theorize them are now available for the next shock, whenever and wherever it arrives.